\renewcommand{\thechapterdisplay}{Trigésimo dia}
\renewcommand{\thechapterrunner}{As últimas vontades dos falecidos}
\chapter[Trigésimo dia --- As últimas vontades dos falecidos]{As últimas vontades dos falecidos}
\section*{É preciso executá-las fielmente}

Se alguma coisa obriga sobre a terra e deve ser respeitada, são os últimos desejos dos moribundos: eles são sagrados. O santo Concílio de Trento recomenda aos bispos que velem atentamente pelo cumprimento dos testamentos piedosos, feitos pelos fiéis falecidos. Outros concílios vão mais longe, a ponto de privar da comunhão eclesial aqueles que se apropriam das doações dos moribundos ou que cumprem de modo diferente as suas últimas vontades. Leis tão severas nos fazem compreender o quanto nos tornamos culpados ao privar os defuntos dos sufrágios que em vida quiseram assegurar para si após sua morte. Portanto, infeliz daquele que engorda-se com o que é devido às almas do Purgatório. Eles as privam do alívio que teriam recebido, constituem-se de algum modo seus algozes, e tornam-se responsáveis por seus sofrimentos diante de Deus. Ah! Se o mundo os absolve desse roubo sacrílego, se o mundo não se escandaliza com esse desprezo, Deus não os absolverá assim facilmente. E o dia virá em que Ele pedirá uma conta rigorosa dessas injustiças que eles nem sequer julgam serem censuráveis. Serão provavelmente punidos, já neste mundo, por castigos temporais, e quem pode calcular quão longas e rigorosas serão as penas que terão de suportar no outro mundo?

Alma cristã, refleti! Vosso pais, vossos amigos, vossos benfeitores, não fizeram eles piedosas recomendações ao morrerem? Não pediram eles com a própria voz ou em testamento, orações, missas, esmolas? Não ao menos suplicaram, com lágrimas, que vos lembrásseis sempre deles diante do Senhor? Pois bem! Fizestes jus à confiança que em vós depositaram? Tendes cumprido plenamente e conscientemente todas as obrigações que eles vos deixaram? Ah! Se vós não as tendes feito, apressai-vos a quitar esse débito sagrado de justiça.


\section*{É preciso executá-las logo}

Não é somente preciso cumprir com fidelidade as últimas vontades dos falecido, mas devemos fazê-lo tão logo quanto possível, afim de não privarmos essas queridas almas do alívio que podem receber quando as atendemos. Seja através das missas que serão celebradas em sua intenção, seja através das orações dos pobres. Sim, quando entregamos aos pobres as esmolas que lhes foram deixadas pelos falecidos, devemos sempre pedir-lhes que rezem por seus benfeitores. Cada dia de atraso é uma falta de que somos responsáveis e que soma ao suplício dos mortos. Ah! Se compreendêssemos como são terríveis os sofrimentos no Purgatório, em vez de mudar ou postergar o cumprimento daquilo que poderia diminuí-los, nós nos apressaríamos em trazer os remédios prontos e eficazes a essas almas benditas, tão dignas de nossa compaixão. Ai de nós! quantos herdeiros inescrupulosos têm graves censuras a fazer a si mesmos, por causa de sua negligência em cumprir os compromissos sagrados que contraíram para com seus pais falecidos. 

Instruí-vos, alma cristã, e não entregais senão a pessoas de confiança o cuidado de executar vossos últimos pedidos. Depositai nas mão mais confiáveis o montante que vós quereis destinar às boas obras ou a que se faça celebrar missas pela vossa libertação, depois que morrerdes. É o meio mais seguro de que seus últimos desejos sejam cumpridos, a menos que tendes a felicidade de pertencer a uma dessas famílias cristãs que, junto com a fé, conservaram o respeito devido à memória dos mortos.


\section*{Exemplo}

Eis aqui um relato que mostra até que ponto são punidos, às vezes, aqueles que não executam as últimas vontades dos moribundos. Conta-se nos Feitos de Carlos Magno que um valente capitão, cuja bravura todos elogiavam, chegara ao fim de sua vida. Ele então chamou um de seus parentes, a quem tinha ajudado muitas vezes, e disse-lhe. \textit{Passei sessenta anos a serviço de meu rei, sem jamais adquirir outra coisa além de meu soldo habitual. Ao morrer, resta-me apenas meu fiel cavalo que me prestou tantos serviços. Quando eu tiver dado o último suspiro, to o venderás e darás o dinheiro aos pobres para o alívio de minha alma.} O parente prometeu. Mas, quando o capitão entregou sua alma a Deus, este homem, seduzido pela beleza e pelas qualidades do animal, guardou-o para si, sem dar aos pobres a esmola combinada. Mal tinha se passado meio ano, quando a alma do defunto apareceu a este parente egoísta, tão infiel à sua promessa e disse-lhe. \textit{Infeliz, tu não executaste minhas últimas vontades, tu não cumpriste teus compromissos. Assim, tu és a causa de todos os tormentos que tenho sofrido, pois minha esmola me teria poupado deles. Pois bem, saiba que tua conduta será punida com uma morte rápida, e que um castigo muito particular te está reservado: tu carregarás a pena devida a tuas próprias faltas, e tu sofrerás em meu lugar todas aquelas que eu deveria ainda sofrer para satisfazer a justiça divina.} O culpado ficou esmagado por esta ameaça e, querendo pôr ordem em sua consciência, apressou-se em cumprir as últimas vontades do defunto, depois fez tudo o que pôde para evitar a morte eterna, mas não pôde evitar a morte do corpo que lhe tinha sido anunciada, e que o levou poucos dias depois. Tanto é verdade que a injustiça e a ingratidão para com os mortos são detestadas por Deus, que Ele as pune neste mundo e no outro.


\textbf{Oremos.} Não permitais, ó meu Deus, que uma indigna ganância ou uma culpável negligência me faça faltar aos meus deveres de justiça para com os mortos. Seus direitos são sagrados, suas últimas vontades serão igualmente sagradas para mim. Satisfarei plenamente a todas as obrigações que me deixarem, feliz se puder, por minha presteza, por minhas orações, apressar a hora de sua libertação. Jesus misericordioso! Maria, rainha do Purgatório! Sede-lhes propícios, e que repousem na paz do Céu!



\begin{center}
\textit{Requiescant in pace!}
\end{center}

Rezai agora uma dezena do terço e as orações no final deste livro.
